\section{Experiments actually performed}

\subsection{Environment and reproducibility}

The execution environment was Linux x86\_64 with Python 3.13.5, NumPy 2.3.5,
SciPy 1.17.0, SymPy 1.14.0, and pytest 9.0.2. The synthetic benchmark uses seed
20260914. Numerical-library thread environment variables were set to one for
the recorded run. Full version strings are in \path{results/benchmarks.json}.
The test run completed with \textbf{73 passed in 1.11 seconds}.

No \code{lean}, \code{lake}, or successful Lean execution route was available.
The supplied Lean comparison runner therefore recorded \code{not\_run}, with
that reason, rather than manufacturing zero failures or baseline timings.
Source inspection, Python execution, and generated Lean source are three
different forms of evidence and remain separately labeled throughout the
package.

The benchmarks were executed once for diagnostic purposes; there are no
repeated-trial confidence intervals or hardware-normalized performance claims.
The exact combinatorial counts and algebraic acceptance results are more
informative here than submillisecond timings.

\subsection{Summary of observed outcomes}

\begin{center}
\begin{tabularx}{\textwidth}{@{}p{0.34\textwidth}rY@{}}
\toprule
Experiment & Outcome & Interpretation\\
\midrule
Pytest suite & 73 passed & Includes parametrized, randomized, negative, and mutation tests\\
Named nonlinear cases & 9 accepted / 12 & Two true cases unresolved; one deliberately false case unresolved\\
Generated product-cone cases & 100 / 100 & Constructed inside the covered cone; not representative problem sampling\\
Separate false-polynomial controls & 0 accepted / 24 & All returned unknown; exact negative witnesses checked separately\\
Bernstein cases & 43 accepted / 45 & One true nondyadic-zero case and one false case unresolved\\
Additive recurrence invariants & 49 / 49 & Polynomial increments; base and universal step identities checked\\
Horn growth family & 5 depths & Both strategies prove all goals; search work differs sharply\\
Actual Lean comparison & Not run & No measured superiority over \grind\\
\bottomrule
\end{tabularx}
\end{center}

The 109 accepted nonlinear certificates are the nine named successes plus the
100 generated cases. They are not 109 previously unsolved Lean theorems. The
49 invariants concern a deliberately restricted polynomial recurrence class,
not arbitrary recursive Lean programs.

\subsection{Demand search: counts before speed claims}

\begin{center}
\begin{tabular}{@{}rrrrr@{}}
\toprule
Depth & Forward facts & Demand atoms & Forward ms & Demand ms\\
\midrule
4 & 31 & 5 & 0.416 & 0.128 \\
8 & 511 & 9 & 6.870 & 0.232 \\
10 & 2,047 & 11 & 35.032 & 0.293 \\
12 & 8,191 & 13 & 153.160 & 0.379 \\
14 & 32,767 & 15 & 846.283 & 0.478 \\
\bottomrule
\end{tabular}
\end{center}


The number of proof nodes in both returned certificates is $d+1$ at every
depth. Thus the search-state reduction is not achieved by silently omitting
necessary proof steps. At depth 14, forward search generated 32,766 rule
instances and demand search generated 14; each yielded the same 15-node chain.
The toy rule order intentionally makes the forward target last at its depth.
A different forward strategy, or the actual \grind\ implementation, can behave
differently.

Correctness testing extends beyond this favorable family. The tests exercise
pure dependency cycles, cycles with a base fact, multiple-premise AND-rules,
wrong goals, invalid proof-parent indices, depth/fact limits, and randomized
finite ground Horn systems compared across both search implementations.
These tests support the implementation's behavior, not a completeness claim
for general theorem proving.

\subsection{Nonlinear cases, including unresolved truths}

The named successful cases are a shifted square, two-variable AM--GM, a cubic
box product, a quartic box product, an equality-ideal identity, a shifted square
with rational margin, two-dimensional Cauchy--Schwarz, a product of three lower
bounds, and three-variable variance. The quartic product under a degree-2 cap
and the Motzkin polynomial remain unresolved. A nearly nonnegative but actually
false shifted square also remains unresolved, as required.

For the generated set, each target is a positive rational sum of six products
of constraints from $x,1-x,y,1-y$, with product degree at most four. This
construction intentionally places the targets inside the chosen certificate
cone. It tests discovery, rational reconstruction, and checking on diverse
coefficient combinations; it does not estimate performance on randomly chosen
mathematical inequalities. The median recorded discovery time was about
21.5 milliseconds and the largest about 103.6 milliseconds for this set.
These are single-run local measurements, not a speed comparison with Lean.

Each of the 24 additional negative controls has the form
$(x-a)^2-\varepsilon$ with $\varepsilon>0$ rational. The benchmark separately
checks its exact negative value at $a$. The cone routine itself returns
\code{unknown}; it does not return a certified refutation. Tests also corrupt
weights, remove sign assumptions, change targets by tiny rational quantities,
and introduce invalid factor indices. The checker must reject those corruptions
even if a numerical optimizer would regard the residual as negligible.

\subsection{Bernstein and invariant evidence}

The Bernstein set contains the worked two-leaf positive polynomial, a box
product, a dyadic zero, a nondyadic zero, a false polynomial, and 40 strictly
positive rational quadratics. The 43 accepted cases have complete checked
subdivision trees. A test corrupts the box-cover geometry to ensure local
positivity does not masquerade as a complete proof.

The invariant set contains nine increments $(n+1)^k$ for $0\leq k\leq8$ and
40 fixed-seed random rational polynomial increments with varied starting
values. Every accepted proposal passes exact base and transition checks.
The median recorded synthesis time was about 4.1 milliseconds. This includes
interpolation and validation but not Lean recursor application, source
elaboration, or kernel checking.

The package exports the nine named cone certificates and all 49 invariant
certificates as JSON. A separate reload-and-validate script accepted all 58
saved records. The additional generated certificates and box trees are
reproducibly rediscovered and checked by the benchmark script; they are not
all serialized as separate proof artifacts.

\subsection{What the Lean export does and does not establish}

The exporter emits nine inequality theorems using explicit nonnegative terms,
rational arithmetic, and polynomial identities. It also emits 49 recurrence
proofs using a generic induction lemma for iteration and the generated exact
base/step expressions. The files are \path{lean/Forge/Certificates.lean} and
\path{lean/Forge/Invariants.lean}.

They contain no deliberate admissions or new axioms. Nevertheless, the source
was not compiled, so it is not a verified Lean deliverable. Elaboration details,
tactic API changes, or dependency integration may require corrections. The
right acceptance criterion is a successful build followed by an axiom and
proof audit, not the plausible appearance of generated source. The included
runner leaves this evidence gap explicit.

\section{A fair evaluation for a real Forge implementation}

\subsection{Benchmark strata and leak prevention}

A subsequent Lean evaluation should use distinct strata: standard \grind\
regressions; nonlinear inequalities; fresh recursive definitions needing
induction; accumulator and fusion lemmas; existential constructions; Boolean
and bit-vector goals; quantified first-order problems; and mixed goals in
which one technique supplies a lemma to another. Keep easy goals, failed goals,
and unsupported goals, not only examples chosen after an extension succeeds.

Partition by theorem family and source module rather than individual nearby
lemmas. Remove the target theorem from any retrieval index. Check for equivalent
lemmas, renamed copies, and specialization leakage. A synthesis benchmark
whose target is already available as a library fact primarily evaluates
retrieval. That can be useful, but it is a different experiment.

Record the complete imported environment, theorem attributes, simplification
sets, seeds, solver versions, and tactic configuration. Reset proof state
between runs; do not let a successful trial add a lemma that helps a later
baseline run.

\subsection{Baselines and ablations}

Compare stock \grind, explicitly configured \grind\ variants, relevant
individual tactics, and a straightforward portfolio. Then compare Forge with
one extension removed at a time: demands, nonlinear cuts, induction
generalization, invariant synthesis, witnesses, or external islands. Compare
the cooperative version with the same engines run independently. Otherwise,
a gain could come entirely from adding a known tactic rather than from the
proposed information exchange.

Use two resource protocols. Under \emph{fixed total resources}, all systems
receive the same wall-clock limit, memory limit, core count, and appropriately
reported heartbeat budget. Under \emph{additive fallback resources}, allow the
unchanged baseline budget plus a declared extension budget, reporting both
budgets and the extra cost. Neither protocol substitutes for the other.

Metrics should include checked solved goals, timeouts, elaboration errors,
unsupported translations, proof bytes, reconstruction time, kernel replay time,
and unexpected axioms. Report paired wins and losses, not only aggregate
success. Repeat timings in randomized order and report confidence intervals
or paired distributions. Include proof checking in the main end-to-end result,
while retaining search-only timing for diagnosis.

The supplied comparison script is much smaller: it tries existing tactics on
nine exported certificate contexts, records compiler diagnostics, and optionally
builds the replay library. It is a smoke-test scaffold, not an implementation
of this full protocol.

\section{Development sequence and acceptance gates}

\textbf{Stage 1: reliable replay.} Compile the generated Lean files, fix any
elaboration issues, and audit their axioms. Implement a typed polynomial
reifier and explicit certificate replay before connecting an external oracle.
Reject malformed certificates and wrong-context assumptions in Lean tests.
A successful stage means that a saved certificate proves the exact intended
Lean theorem, not that discovery is fast.

\textbf{Stage 2: nonlinear integration.} Connect the bounded LP proposal service
to the reifier and checker. Add the goal-sensitive square dictionary, product
cuts, and failure diagnostics. Start with real polynomial inequalities and
explicit hypotheses; postpone complicated casts and rational expressions until
their transport lemmas are tested. Measure against \code{nlinarith} and stock
\grind, including cases they already solve.

\textbf{Stage 3: structural coordination.} Implement the scoped obligation graph,
tabled demands, and one-variable structural induction with changed-argument
generalization. Port additive invariant synthesis and prove its recurrence
transport. Demonstrate mixed examples in which synthesized structure feeds
facts to the algebraic engine. Keep a replay script for every success.

\textbf{Stage 4: richer synthesis and delegation.} Add typed witness grammars,
selected solver islands, and optional broader certificate search. Each backend
must pass a proof-rule coverage matrix and strict/native trust-profile audit.
A missing or incompatible backend should disable only its lane.

\textbf{Stage 5: validated deployment.} Run the controlled repository-scale
benchmarks, tune scheduling on a separate training set, preserve stock behavior
in conservative mode, and publish wins, regressions, unsupported cases, and
proof-checking costs. Do not release an experimental always-on policy merely
because it improves an aggregate score on its tuning examples.

These are dependency-ordered milestones, not a promised schedule. The highest
risk is the typed, scoped, version-compatible bridge into Lean and the quality
of the cooperative search policy, not the small algebraic checker alone.

\section{Conclusion}

The credible path beyond \grind\ is to supply the proof structures and
certificates that local saturation lacks, while retaining its mature core.
Forge makes that path concrete: demand-directed obligations, exact nonlinear
certificates, complete box proofs, justified induction generalization,
polynomial invariants, typed witnesses, and narrowly translated solver islands.
Its organizing rule is that speculative search may be aggressive but every
logical consequence must have a checkable derivation in its actual scope.

The executable work supports four parts of this proposal and exposes meaningful
limitations. It does not establish a finished Lean tactic or measured dominance
over \grind. The next decisive evidence is a typed Lean implementation with
successful kernel replay and a paired, controlled evaluation. The supplied
source, saved certificates, tests, and integration contracts make that a
specific engineering program rather than a vague aspiration.
